
\documentclass{article}
\usepackage{colortbl}
\usepackage{makecell}
\usepackage{multirow}
\usepackage{supertabular}

\begin{document}

\newcounter{utterance}

\centering \large Interaction Transcript for game `hot\_air\_balloon', experiment `air\_balloon\_survival\_de\_complexity\_easy', episode 1 with deepseek{-}chat{-}v3.1{-}t1.0.
\vspace{24pt}

{ \footnotesize  \setcounter{utterance}{1}
\setlength{\tabcolsep}{0pt}
\begin{supertabular}{c@{$\;$}|p{.15\linewidth}@{}p{.15\linewidth}p{.15\linewidth}p{.15\linewidth}p{.15\linewidth}p{.15\linewidth}}
    \# & $\;$A & \multicolumn{4}{c}{Game Master} & $\;\:$B\\
    \hline

    \theutterance \stepcounter{utterance}  
    & & \multicolumn{4}{p{0.6\linewidth}}{
        \cellcolor[rgb]{0.9,0.9,0.9}{
            \makecell[{{p{\linewidth}}}]{
                \texttt{\tiny{[P1$\langle$GM]}}
                \texttt{Du nimmst an einem kollaborativen Verhandlungspiel Teil.} \\
\\ 
\texttt{Zusammen mit einem anderen Teilnehmer musst du dich auf eine Reihe von Gegenständen entscheiden, die behalten werden. Jeder von euch hat eine persönliche Verteilung über die Wichtigkeit der einzelnen Gegenstände. Jeder von euch hat eine eigene Meinung darüber, wie wichtig jeder einzelne Gegenstand ist (Gegenstandswichtigkeit). Du kennst die Wichtigkeitsverteilung des anderen Spielers nicht. Zusätzlich siehst du, wie viel Aufwand jeder Gegenstand verursacht.  } \\
\texttt{Ihr dürft euch nur auf eine Reihe von Gegenständen einigen, wenn der Gesamtaufwand der ausgewählten Gegenstände den Maximalaufwand nicht überschreitet:} \\
\\ 
\texttt{Maximalaufwand = 8170} \\
\\ 
\texttt{Hier sind die einzelnen Aufwände der Gegenstände:} \\
\\ 
\texttt{Aufwand der Gegenstände = \{"C32": 594, "C52": 551, "B54": 634, "C16": 204, "C50": 557, "A08": 440, "C59": 734, "B51": 893, "B55": 246, "B59": 592, "A77": 143, "C14": 568, "C24": 472, "B24": 401, "C31": 729, "A52": 201, "C61": 85, "B79": 641, "C54": 958, "B56": 80, "A40": 157, "B11": 806, "B60": 684, "A81": 59, "A95": 31, "B38": 765, "C51": 416, "C03": 392, "A91": 428, "A93": 699, "C28": 141, "C05": 606, "B75": 612, "B96": 133, "B86": 689\}} \\
\\ 
\texttt{Hier ist deine persönliche Verteilung der Wichtigkeit der einzelnen Gegenstände:} \\
\\ 
\texttt{Werte der Gegenstandswichtigkeit = \{"C32": 138, "C52": 583, "B54": 868, "C16": 822, "C50": 783, "A08": 65, "C59": 262, "B51": 121, "B55": 508, "B59": 780, "A77": 461, "C14": 484, "C24": 668, "B24": 389, "C31": 808, "A52": 215, "C61": 97, "B79": 500, "C54": 30, "B56": 915, "A40": 856, "B11": 400, "B60": 444, "A81": 623, "A95": 781, "B38": 786, "C51": 3, "C03": 713, "A91": 457, "A93": 273, "C28": 739, "C05": 822, "B75": 235, "B96": 606, "B86": 105\}} \\
\\ 
\texttt{Ziel:} \\
\\ 
\texttt{Dein Ziel ist es, eine Reihe von Gegenständen auszuhandeln, die dir möglichst viel bringt (d. h. Gegenständen, die DEINE Wichtigkeit maximieren), wobei der Maximalaufwand eingehalten werden muss. Du musst nicht in jeder Nachricht einen VORSCHLAG machen – du kannst auch nur verhandeln. Alle Taktiken sind erlaubt!} \\
\\ 
\texttt{Interaktionsprotokoll:} \\
\\ 
\texttt{Du darfst nur die folgenden strukturierten Formate in deinen Nachrichten verwenden:} \\
\\ 
\texttt{VORSCHLAG: \{'A', 'B', 'C', …\}} \\
\texttt{Schlage einen Deal mit genau diesen Gegenstände vor.} \\
\texttt{ABLEHNUNG: \{'A', 'B', 'C', …\}} \\
\texttt{Lehne den Vorschlag des Gegenspielers ausdrücklich ab.} \\
\texttt{ARGUMENT: \{'...'\}} \\
\texttt{Verteidige deinen letzten Vorschlag oder argumentiere gegen den Vorschlag des Gegenspielers.} \\
\texttt{ZUSTIMMUNG: \{'A', 'B', 'C', …\}} \\
\texttt{Akzeptiere den Vorschlag des Gegenspielers, wodurch das Spiel endet.} \\
\texttt{STRATEGISCHE ÜBERLEGUNGEN: \{'...'\}} \\
\texttt{	Beschreibe strategische Überlegungen, die deine nächsten Schritte erklären. Dies ist eine versteckte Nachricht, die nicht mit dem anderen Teilnehmer geteilt wird.} \\
\\ 
\texttt{Regeln:} \\
\\ 
\texttt{Du darst nur einen Vorschlag mit ZUSTIMMUNG akzeptieren, der vom anderen Spieler zuvor mit VORSCHLAG eingebracht wurde.} \\
\texttt{Du darfst nur Vorschläge mit ABLEHNUNG ablehnen, die vom anderen Spieler durch VORSCHLAG zuvor genannt wurden. } \\
\texttt{Der Gesamtaufwand einer VORSCHLAG{-} oder ZUSTIMMUNG{-}Menge darf nicht größer als der Maximalaufwand sein.  } \\
\texttt{Offenbare deine versteckte Wichtigkeitsverteilung nicht.} \\
\texttt{Ein Schlagwort muss gemäß der Formatvorgaben von einem Doppelpunkt und einem Leerzeichen gefolgt sein. Das Argument muss eine Python{-}Menge mit 0 oder mehr Strings sein.  } \\
\texttt{Die Form muss also SCHLAGWORT: \{...\} sein.} \\
\texttt{Folge nur dem Interaktionsprotokoll und schreibe NICHTS außerhalb der Formatvorgaben.} \\
\texttt{Das Spiel endet, wenn eine Seite mit ZUSTIMMUNG auf einen VORSCHLAG der anderen Seite reagiert.  } \\
\texttt{Der Teil deiner Antwort, der dem anderen Spieler übergeben werden kann, darf nicht leer sein.  } \\
\texttt{Nur Vorschläge, die mit dem Format VORSCHLAG übermittelt wurden und nicht mit ABLEHNUNG abgelehnt wurden, gelten als aktiv.  } \\
\texttt{Jeder deiner Nachrichten muss an mindestens einer Stelle das ARGUMENT Format beinhalten, um in der Verhandlung voranzukommen.} \\
\texttt{Du musst exakt ein mal zu Beginn jeder deiner Nachrichten das 'STRATEGISCHE ÜBERLEGUNGEN' Format verwenden und nicht häufiger. Der Inhalt wird dem anderen Spieler nicht übermittelt, daher kannst du alles einfügen, was du möchtest {-} einschließlich deiner eigenen Wichtigkeitsverteilung. Hier solltest du Schritt für Schritt überlegen, um deinen nächsten Zug zu planen.} \\
            }
        }
    }
    & \\ \\

    \theutterance \stepcounter{utterance}  
    & \multicolumn{4}{p{0.6\linewidth}}{
        \cellcolor[rgb]{0.9,0.9,0.9}{
            \makecell[{{p{\linewidth}}}]{
                \texttt{\tiny{[P1$\rangle$GM]}}
                \texttt{STRATEGISCHE ÜBERLEGUNGEN: \{'Ich muss eine Strategie entwickeln, die meine Gesamtwichtigkeit maximiert, während der Gesamtaufwand unter 8170 bleibt. Meine wichtigsten Gegenstände sind B56 (915), B54 (868), A40 (856), C16 (822), C05 (822), C31 (808), B38 (786), A95 (781), B59 (780), C50 (783) mit hoher Wichtigkeit und niedrigem bis moderatem Aufwand. Niedrige Wichtigkeit bei hohem Aufwand wie C54 (30, 958) sollte ich vermeiden. Ich sollte versuchen, viele hochwertige, niedrigaufwändige Gegenstände zu bekommen und dem Gegner hochaufwändige, für mich unwichtige Gegenstände überlassen.'\}} \\
\\ 
\texttt{ARGUMENT: \{'Ich schlage einen ausgewogenen Deal vor, der für beide Seiten vorteilhaft sein sollte. Mein Fokus liegt auf Gegenständen mit hohem Wert für mich, während ich bereit bin, bei anderen Kompromisse einzugehen.'\}} \\
\\ 
\texttt{VORSCHLAG: \{'B56', 'B54', 'A40', 'C16', 'C05', 'C31', 'B38', 'A95', 'B59', 'C50', 'C52', 'B55', 'C24', 'A81', 'C28', 'B96', 'C03', 'A91', 'B60', 'A77', 'B79', 'C32', 'C61', 'A52', 'B75', 'B86', 'A93'\}} \\
            }
        }
    }
    & & \\ \\

    \theutterance \stepcounter{utterance}  
    & & \multicolumn{4}{p{0.6\linewidth}}{
        \cellcolor[rgb]{0.9,0.9,0.9}{
            \makecell[{{p{\linewidth}}}]{
                \texttt{\tiny{[GM$\rangle$P2]}}
                \texttt{Du nimmst an einem kollaborativen Verhandlungspiel Teil.} \\
\\ 
\texttt{Zusammen mit einem anderen Teilnehmer musst du dich auf eine Reihe von Gegenständen entscheiden, die behalten werden. Jeder von euch hat eine persönliche Verteilung über die Wichtigkeit der einzelnen Gegenstände. Jeder von euch hat eine eigene Meinung darüber, wie wichtig jeder einzelne Gegenstand ist (Gegenstandswichtigkeit). Du kennst die Wichtigkeitsverteilung des anderen Spielers nicht. Zusätzlich siehst du, wie viel Aufwand jeder Gegenstand verursacht.  } \\
\texttt{Ihr dürft euch nur auf eine Reihe von Gegenständen einigen, wenn der Gesamtaufwand der ausgewählten Gegenstände den Maximalaufwand nicht überschreitet:} \\
\\ 
\texttt{Maximalaufwand = 8170} \\
\\ 
\texttt{Hier sind die einzelnen Aufwände der Gegenstände:} \\
\\ 
\texttt{Aufwand der Gegenstände = \{"C32": 594, "C52": 551, "B54": 634, "C16": 204, "C50": 557, "A08": 440, "C59": 734, "B51": 893, "B55": 246, "B59": 592, "A77": 143, "C14": 568, "C24": 472, "B24": 401, "C31": 729, "A52": 201, "C61": 85, "B79": 641, "C54": 958, "B56": 80, "A40": 157, "B11": 806, "B60": 684, "A81": 59, "A95": 31, "B38": 765, "C51": 416, "C03": 392, "A91": 428, "A93": 699, "C28": 141, "C05": 606, "B75": 612, "B96": 133, "B86": 689\}} \\
\\ 
\texttt{Hier ist deine persönliche Verteilung der Wichtigkeit der einzelnen Gegenstände:} \\
\\ 
\texttt{Werte der Gegenstandswichtigkeit = \{"C32": 138, "C52": 583, "B54": 868, "C16": 822, "C50": 783, "A08": 65, "C59": 262, "B51": 121, "B55": 508, "B59": 780, "A77": 461, "C14": 484, "C24": 668, "B24": 389, "C31": 808, "A52": 215, "C61": 97, "B79": 500, "C54": 30, "B56": 915, "A40": 856, "B11": 400, "B60": 444, "A81": 623, "A95": 781, "B38": 786, "C51": 3, "C03": 713, "A91": 457, "A93": 273, "C28": 739, "C05": 822, "B75": 235, "B96": 606, "B86": 105\}} \\
\\ 
\texttt{Ziel:} \\
\\ 
\texttt{Dein Ziel ist es, eine Reihe von Gegenständen auszuhandeln, die dir möglichst viel bringt (d. h. Gegenständen, die DEINE Wichtigkeit maximieren), wobei der Maximalaufwand eingehalten werden muss. Du musst nicht in jeder Nachricht einen VORSCHLAG machen – du kannst auch nur verhandeln. Alle Taktiken sind erlaubt!} \\
\\ 
\texttt{Interaktionsprotokoll:} \\
\\ 
\texttt{Du darfst nur die folgenden strukturierten Formate in deinen Nachrichten verwenden:} \\
\\ 
\texttt{VORSCHLAG: \{'A', 'B', 'C', …\}} \\
\texttt{Schlage einen Deal mit genau diesen Gegenstände vor.} \\
\texttt{ABLEHNUNG: \{'A', 'B', 'C', …\}} \\
\texttt{Lehne den Vorschlag des Gegenspielers ausdrücklich ab.} \\
\texttt{ARGUMENT: \{'...'\}} \\
\texttt{Verteidige deinen letzten Vorschlag oder argumentiere gegen den Vorschlag des Gegenspielers.} \\
\texttt{ZUSTIMMUNG: \{'A', 'B', 'C', …\}} \\
\texttt{Akzeptiere den Vorschlag des Gegenspielers, wodurch das Spiel endet.} \\
\texttt{STRATEGISCHE ÜBERLEGUNGEN: \{'...'\}} \\
\texttt{	Beschreibe strategische Überlegungen, die deine nächsten Schritte erklären. Dies ist eine versteckte Nachricht, die nicht mit dem anderen Teilnehmer geteilt wird.} \\
\\ 
\texttt{Regeln:} \\
\\ 
\texttt{Du darst nur einen Vorschlag mit ZUSTIMMUNG akzeptieren, der vom anderen Spieler zuvor mit VORSCHLAG eingebracht wurde.} \\
\texttt{Du darfst nur Vorschläge mit ABLEHNUNG ablehnen, die vom anderen Spieler durch VORSCHLAG zuvor genannt wurden. } \\
\texttt{Der Gesamtaufwand einer VORSCHLAG{-} oder ZUSTIMMUNG{-}Menge darf nicht größer als der Maximalaufwand sein.  } \\
\texttt{Offenbare deine versteckte Wichtigkeitsverteilung nicht.} \\
\texttt{Ein Schlagwort muss gemäß der Formatvorgaben von einem Doppelpunkt und einem Leerzeichen gefolgt sein. Das Argument muss eine Python{-}Menge mit 0 oder mehr Strings sein.  } \\
\texttt{Die Form muss also SCHLAGWORT: \{...\} sein.} \\
\texttt{Folge nur dem Interaktionsprotokoll und schreibe NICHTS außerhalb der Formatvorgaben.} \\
\texttt{Das Spiel endet, wenn eine Seite mit ZUSTIMMUNG auf einen VORSCHLAG der anderen Seite reagiert.  } \\
\texttt{Der Teil deiner Antwort, der dem anderen Spieler übergeben werden kann, darf nicht leer sein.  } \\
\texttt{Nur Vorschläge, die mit dem Format VORSCHLAG übermittelt wurden und nicht mit ABLEHNUNG abgelehnt wurden, gelten als aktiv.  } \\
\texttt{Jeder deiner Nachrichten muss an mindestens einer Stelle das ARGUMENT Format beinhalten, um in der Verhandlung voranzukommen.} \\
\texttt{Du musst exakt ein mal zu Beginn jeder deiner Nachrichten das 'STRATEGISCHE ÜBERLEGUNGEN' Format verwenden und nicht häufiger. Der Inhalt wird dem anderen Spieler nicht übermittelt, daher kannst du alles einfügen, was du möchtest {-} einschließlich deiner eigenen Wichtigkeitsverteilung. Hier solltest du Schritt für Schritt überlegen, um deinen nächsten Zug zu planen.} \\
\\ 
\texttt{Du erhältst nun die erste Nachricht des anderen Spielers.} \\
\\ 
\texttt{ARGUMENT: \{'Ich schlage einen ausgewogenen Deal vor, der für beide Seiten vorteilhaft sein sollte. Mein Fokus liegt auf Gegenständen mit hohem Wert für mich, während ich bereit bin, bei anderen Kompromisse einzugehen.'\}} \\
\\ 
\texttt{VORSCHLAG: \{'B56', 'B54', 'A40', 'C16', 'C05', 'C31', 'B38', 'A95', 'B59', 'C50', 'C52', 'B55', 'C24', 'A81', 'C28', 'B96', 'C03', 'A91', 'B60', 'A77', 'B79', 'C32', 'C61', 'A52', 'B75', 'B86', 'A93'\}} \\
            }
        }
    }
    & \\ \\

    \theutterance \stepcounter{utterance}  
    & & & \multicolumn{4}{p{0.6\linewidth}}{
        \cellcolor[rgb]{0.9,0.9,0.9}{
            \makecell[{{p{\linewidth}}}]{
                \texttt{\tiny{[GM$\langle$P2]}}
                \texttt{STRATEGISCHE ÜBERLEGUNGEN: \{'Der Gegner hat einen Vorschlag mit 27 Gegenständen gemacht. Ich muss prüfen, ob der Gesamtaufwand unter 8170 liegt und wie viel Wert dieser Vorschlag für mich hat. Basierend auf meiner Wichtigkeitsverteilung berechne ich den Gesamtwert und vergleiche mit möglichen Alternativen. Ich sollte Gegenstände mit hohem Wert für mich (wie B56:915, B54:868, C16:822, C05:822, C31:808, B38:786, A95:781, B59:780, C50:783, A40:856) behalten, während ich niedrige Werte (wie C51:3, C54:30, A08:65, C61:97, B86:105, B51:121) opfern kann. Der Vorschlag des Gegners enthält viele meiner High{-}Value{-}Items, aber auch einige niedrige. Ich sollte argumentieren und einen counter Vorschlag machen, der meinen Wert optimiert.'\}} \\
\\ 
\texttt{ARGUMENT: \{'Vielen Dank für Ihren Vorschlag. Während ich einige der enthaltenen Gegenstände schätze, gibt es bestimmte High{-}Value{-}Items für mich, die fehlen, und einige mit geringem Wert für mich, die ich lieber entfernen würde. Lassen Sie uns einen Deal finden, der für beide Seiten optimal ist.'\}} \\
\\ 
\texttt{VORSCHLAG: \{'B56', 'B54', 'A40', 'C16', 'C05', 'C31', 'B38', 'A95', 'B59', 'C50', 'C52', 'B55', 'C24', 'A81', 'C28', 'B96', 'C03', 'A91', 'B60', 'A77', 'B79', 'C32', 'C61', 'A52', 'B75', 'B86', 'A93', 'C14', 'A52'\}} \\
            }
        }
    }
     \\ \\

    \theutterance \stepcounter{utterance}  
    & & \multicolumn{4}{p{0.6\linewidth}}{
        \cellcolor[rgb]{0.9,0.9,0.9}{
            \makecell[{{p{\linewidth}}}]{
                \texttt{\tiny{[P1$\langle$GM]}}
                \texttt{ARGUMENT: \{'Vielen Dank für Ihren Vorschlag. Während ich einige der enthaltenen Gegenstände schätze, gibt es bestimmte High{-}Value{-}Items für mich, die fehlen, und einige mit geringem Wert für mich, die ich lieber entfernen würde. Lassen Sie uns einen Deal finden, der für beide Seiten optimal ist.'\}} \\
\\ 
\texttt{VORSCHLAG: \{'B56', 'B54', 'A40', 'C16', 'C05', 'C31', 'B38', 'A95', 'B59', 'C50', 'C52', 'B55', 'C24', 'A81', 'C28', 'B96', 'C03', 'A91', 'B60', 'A77', 'B79', 'C32', 'C61', 'A52', 'B75', 'B86', 'A93', 'C14', 'A52'\}} \\
            }
        }
    }
    & \\ \\

    \theutterance \stepcounter{utterance}  
    & \multicolumn{4}{p{0.6\linewidth}}{
        \cellcolor[rgb]{0.9,0.9,0.9}{
            \makecell[{{p{\linewidth}}}]{
                \texttt{\tiny{[P1$\rangle$GM]}}
                \texttt{STRATEGISCHE ÜBERLEGUNGEN: \{'Der Gegner hat meinen Vorschlag modifiziert und A52 doppelt eingefügt (was unmöglich ist) und C14 hinzugefügt. Ich muss prüfen, ob dieser Vorschlag für mich vorteilhaft ist. Der Gesamtaufwand beträgt 594+551+634+204+557+143+734+246+592+568+472+401+729+201+85+641+80+157+806+684+59+31+765+416+392+428+699+141+606+612+133+689 = 13239, was den Maximalaufwand von 8170 deutlich überschreitet. Dies ist ein ungültiger Vorschlag. Ich sollte ihn ablehnen und einen counter{-}Vorschlag machen, der unter dem Budget bleibt und meine wichtigsten Items enthält.'\}} \\
\\ 
\texttt{ARGUMENT: \{'Ihr Vorschlag überschreitet den maximal erlaubten Aufwand von 8170 deutlich. Zudem ist Gegenstand A52 doppelt aufgeführt. Ich schlage einen realistischen Counter{-}Vorschlag vor, der innerhalb des Budgets bleibt und für beide Seiten vorteilhaft ist.'\}} \\
\\ 
\texttt{ABLEHNUNG: \{'B56', 'B54', 'A40', 'C16', 'C05', 'C31', 'B38', 'A95', 'B59', 'C50', 'C52', 'B55', 'C24', 'A81', 'C28', 'B96', 'C03', 'A91', 'B60', 'A77', 'B79', 'C32', 'C61', 'A52', 'B75', 'B86', 'A93', 'C14', 'A52'\}} \\
\\ 
\texttt{VORSCHLAG: \{'B56', 'B54', 'A40', 'C16', 'C05', 'C31', 'B38', 'A95', 'B59', 'C50', 'C52', 'C24', 'A81', 'C28', 'B96', 'C03', 'A91', 'B60', 'A77', 'B79', 'C61', 'B75', 'B86', 'C32', 'B55'\}} \\
            }
        }
    }
    & & \\ \\

    \theutterance \stepcounter{utterance}  
    & & \multicolumn{4}{p{0.6\linewidth}}{
        \cellcolor[rgb]{0.9,0.9,0.9}{
            \makecell[{{p{\linewidth}}}]{
                \texttt{\tiny{[GM$\rangle$P2]}}
                \texttt{ARGUMENT: \{'Ihr Vorschlag überschreitet den maximal erlaubten Aufwand von 8170 deutlich. Zudem ist Gegenstand A52 doppelt aufgeführt. Ich schlage einen realistischen Counter{-}Vorschlag vor, der innerhalb des Budgets bleibt und für beide Seiten vorteilhaft ist.'\}} \\
\\ 
\texttt{ABLEHNUNG: \{'B56', 'B54', 'A40', 'C16', 'C05', 'C31', 'B38', 'A95', 'B59', 'C50', 'C52', 'B55', 'C24', 'A81', 'C28', 'B96', 'C03', 'A91', 'B60', 'A77', 'B79', 'C32', 'C61', 'A52', 'B75', 'B86', 'A93', 'C14', 'A52'\}} \\
\\ 
\texttt{VORSCHLAG: \{'B56', 'B54', 'A40', 'C16', 'C05', 'C31', 'B38', 'A95', 'B59', 'C50', 'C52', 'C24', 'A81', 'C28', 'B96', 'C03', 'A91', 'B60', 'A77', 'B79', 'C61', 'B75', 'B86', 'C32', 'B55'\}} \\
            }
        }
    }
    & \\ \\

    \theutterance \stepcounter{utterance}  
    & & & \multicolumn{4}{p{0.6\linewidth}}{
        \cellcolor[rgb]{0.9,0.9,0.9}{
            \makecell[{{p{\linewidth}}}]{
                \texttt{\tiny{[GM$\langle$P2]}}
                \texttt{STRATEGISCHE ÜBERLEGUNGEN: \{'Der Gegner hat meinen vorherigen Vorschlag abgelehnt weil er den Maximalaufwand überschreitet (ich hatte A52 doppelt und C14 hinzugefügt). Sein neuer Vorschlag entfernt B55 und behält 25 Gegenstände. Ich muss den Aufwand und Wert berechnen: Aufwand = 80+634+157+204+606+729+765+31+780+557+551+472+59+141+133+392+428+684+143+641+85+612+689+594+508? = 80+634=714, +157=871, +204=1075, +606=1681, +729=2410, +765=3175, +31=3206, +780=3986, +557=4543, +551=5094, +472=5566, +59=5625, +141=5766, +133=5899, +392=6291, +428=6719, +684=7403, +143=7546, +641=8187, +85=8272, +612=8884, +689=9573, +594=10167, +508=10675 {-} das überschreitet 8170 bei weitem! Ich muss den Aufwand korrekt berechnen. Vielleicht hat er B55 entfernt (246 Aufwand) aber die Liste ist unklar. Besser ich berechne systematisch mit den gegebenen Aufwänden.'\}} \\
\\ 
\texttt{ARGUMENT: \{'Ihr Vorschlag scheint den Maximalaufwand von 8170 immer noch deutlich zu überschreiten. Lassen Sie uns einen realistischen Vorschlag machen, der tatsächlich innerhalb des Budgets liegt. Ich schlage vor, wir konzentrieren uns auf die Gegenstände mit dem besten Wert{-}Aufwand{-}Verhältnis.'\}} \\
\\ 
\texttt{ABLEHNUNG: \{'B56', 'B54', 'A40', 'C16', 'C05', 'C31', 'B38', 'A95', 'B59', 'C50', 'C52', 'C24', 'A81', 'C28', 'B96', 'C03', 'A91', 'B60', 'A77', 'B79', 'C61', 'B75', 'B86', 'C32', 'B55'\}} \\
\\ 
\texttt{VORSCHLAG: \{'B56', 'B54', 'A40', 'C16', 'C05', 'C31', 'B38', 'A95', 'B59', 'C50', 'C24', 'A81', 'C28', 'B96', 'C03', 'B60', 'A77', 'B79', 'C61', 'B75', 'B86', 'C32'\}} \\
            }
        }
    }
     \\ \\

    \theutterance \stepcounter{utterance}  
    & & \multicolumn{4}{p{0.6\linewidth}}{
        \cellcolor[rgb]{0.9,0.9,0.9}{
            \makecell[{{p{\linewidth}}}]{
                \texttt{\tiny{[P1$\langle$GM]}}
                \texttt{ARGUMENT: \{'Ihr Vorschlag scheint den Maximalaufwand von 8170 immer noch deutlich zu überschreiten. Lassen Sie uns einen realistischen Vorschlag machen, der tatsächlich innerhalb des Budgets liegt. Ich schlage vor, wir konzentrieren uns auf die Gegenstände mit dem besten Wert{-}Aufwand{-}Verhältnis.'\}} \\
\\ 
\texttt{ABLEHNUNG: \{'B56', 'B54', 'A40', 'C16', 'C05', 'C31', 'B38', 'A95', 'B59', 'C50', 'C52', 'C24', 'A81', 'C28', 'B96', 'C03', 'A91', 'B60', 'A77', 'B79', 'C61', 'B75', 'B86', 'C32', 'B55'\}} \\
\\ 
\texttt{VORSCHLAG: \{'B56', 'B54', 'A40', 'C16', 'C05', 'C31', 'B38', 'A95', 'B59', 'C50', 'C24', 'A81', 'C28', 'B96', 'C03', 'B60', 'A77', 'B79', 'C61', 'B75', 'B86', 'C32'\}} \\
            }
        }
    }
    & \\ \\

    \theutterance \stepcounter{utterance}  
    & \multicolumn{4}{p{0.6\linewidth}}{
        \cellcolor[rgb]{0.9,0.9,0.9}{
            \makecell[{{p{\linewidth}}}]{
                \texttt{\tiny{[P1$\rangle$GM]}}
                \texttt{STRATEGISCHE ÜBERLEGUNGEN: \{'Der Gegner hat recht {-} mein letzter Vorschlag war immer noch über dem Budget. Ich muss den Aufwand berechnen: B56(80) + B54(634) + A40(157) + C16(204) + C05(606) + C31(729) + B38(765) + A95(31) + B59(592) + C50(557) + C24(472) + A81(59) + C28(141) + B96(133) + C03(392) + B60(684) + A77(143) + B79(641) + C61(85) + B75(612) + B86(689) + C32(594) = 80+634+157+204+606+729+765+31+592+557+472+59+141+133+392+684+143+641+85+612+689+594 = 10541, immer noch über 8170. Ich muss die Liste reduzieren und mich auf meine wichtigsten Items konzentrieren.'\}} \\
\\ 
\texttt{ARGUMENT: \{'Sie haben recht, mein Vorschlag war nicht budgetkonform. Ich konzentriere mich jetzt auf die Gegenstände mit dem besten Wert{-}Aufwand{-}Verhältnis für mich und stelle sicher, dass der Gesamtaufwand unter 8170 bleibt.'\}} \\
\\ 
\texttt{ABLEHNUNG: \{'B56', 'B54', 'A40', 'C16', 'C05', 'C31', 'B38', 'A95', 'B59', 'C50', 'C24', 'A81', 'C28', 'B96', 'C03', 'B60', 'A77', 'B79', 'C61', 'B75', 'B86', 'C32'\}} \\
\\ 
\texttt{VORSCHLAG: \{'B56', 'B54', 'A40', 'C16', 'C05', 'C31', 'B38', 'A95', 'B59', 'C50', 'C24', 'A81', 'C28', 'B96', 'C03', 'A77', 'C61', 'B75', 'C32'\}} \\
            }
        }
    }
    & & \\ \\

    \theutterance \stepcounter{utterance}  
    & & \multicolumn{4}{p{0.6\linewidth}}{
        \cellcolor[rgb]{0.9,0.9,0.9}{
            \makecell[{{p{\linewidth}}}]{
                \texttt{\tiny{[GM$\rangle$P2]}}
                \texttt{ARGUMENT: \{'Sie haben recht, mein Vorschlag war nicht budgetkonform. Ich konzentriere mich jetzt auf die Gegenstände mit dem besten Wert{-}Aufwand{-}Verhältnis für mich und stelle sicher, dass der Gesamtaufwand unter 8170 bleibt.'\}} \\
\\ 
\texttt{ABLEHNUNG: \{'B56', 'B54', 'A40', 'C16', 'C05', 'C31', 'B38', 'A95', 'B59', 'C50', 'C24', 'A81', 'C28', 'B96', 'C03', 'B60', 'A77', 'B79', 'C61', 'B75', 'B86', 'C32'\}} \\
\\ 
\texttt{VORSCHLAG: \{'B56', 'B54', 'A40', 'C16', 'C05', 'C31', 'B38', 'A95', 'B59', 'C50', 'C24', 'A81', 'C28', 'B96', 'C03', 'A77', 'C61', 'B75', 'C32'\}} \\
            }
        }
    }
    & \\ \\

    \theutterance \stepcounter{utterance}  
    & & & \multicolumn{4}{p{0.6\linewidth}}{
        \cellcolor[rgb]{0.9,0.9,0.9}{
            \makecell[{{p{\linewidth}}}]{
                \texttt{\tiny{[GM$\langle$P2]}}
                \texttt{STRATEGISCHE ÜBERLEGUNGEN: \{'Der Gegner schlägt 19 Gegenstände vor. Ich muss den Gesamtaufwand berechnen: B56:80, B54:634, A40:157, C16:204, C05:606, C31:729, B38:765, A95:31, B59:592, C50:557, C24:472, A81:59, C28:141, B96:133, C03:392, A77:143, C61:85, B75:612, C32:594. Summe = 80+634=714, +157=871, +204=1075, +606=1681, +729=2410, +765=3175, +31=3206, +592=3798, +557=4355, +472=4827, +59=4886, +141=5027, +133=5160, +392=5552, +143=5695, +85=5780, +612=6392, +594=6986. Dies ist unter 8170. Jetzt berechne ich den Wert für mich: B56:915, B54:868, A40:856, C16:822, C05:822, C31:808, B38:786, A95:781, B59:780, C50:783, C24:668, A81:623, C28:739, B96:606, C03:713, A77:461, C61:97, B75:235, C32:138. Summe = 915+868=1783, +856=2639, +822=3461, +822=4283, +808=5091, +786=5877, +781=6658, +780=7438, +783=8221, +668=8889, +623=9512, +739=10251, +606=10857, +713=11570, +461=12031, +97=12128, +235=12363, +138=12501. Das ist ein sehr guter Wert für mich! Allerdings enthält der Vorschlag einige niedrigwertige Items wie C61:97 und B75:235. Vielleicht kann ich diese entfernen und durch höherwertige ersetzen, aber der Vorschlag ist bereits sehr gut. Ich sollte zustimmen.'\}} \\
\\ 
\texttt{ARGUMENT: \{'Dieser Vorschlag ist deutlich besser budgetkonform und enthält viele wertvolle Gegenstände für mich. Ich bin bereit, diesen Deal zu akzeptieren.'\}} \\
\\ 
\texttt{ZUSTIMMUNG: \{'B56', 'B54', 'A40', 'C16', 'C05', 'C31', 'B38', 'A95', 'B59', 'C50', 'C24', 'A81', 'C28', 'B96', 'C03', 'A77', 'C61', 'B75', 'C32'\}} \\
            }
        }
    }
     \\ \\

    \theutterance \stepcounter{utterance}  
    & & & \multicolumn{2}{p{0.3\linewidth}}{
        \cellcolor[rgb]{0.9,0.9,0.9}{
            \makecell[{{p{\linewidth}}}]{
                \texttt{\tiny{[GM$|$GM]}}
                \texttt{game successful} \\
            }
        }
    }
    & & \\ \\

    \theutterance \stepcounter{utterance}  
    & & & \multicolumn{2}{p{0.3\linewidth}}{
        \cellcolor[rgb]{0.9,0.9,0.9}{
            \makecell[{{p{\linewidth}}}]{
                \texttt{\tiny{[GM$|$GM]}}
                \texttt{end game} \\
            }
        }
    }
    & & \\ \\

\end{supertabular}
}

\end{document}
